\chapter{Delirium}

\section*{Definition}
Acute confusional state characterized by fluctuating consciousness impaired attention disorientation altered perception cognitive dysfunction.

\section*{Pathophysiology}
Medical emergency requiring rapid identification of underlying cause; differs from chronic confusion/dementia by abrupt onset fluctuating course reversible etiations hyperactive hypoactive mixed subtypes associated with increased mortality morbidity prolonged hospitalization long-term cognitive decline especially elderly population.

\section*{Etiology}
\begin{itemize}
  \item Infection/sepsis
  \item Medication toxicity/polypharmacy
  \item Metabolic abnormalities
  \item Alcohol withdrawal
  \item Postoperative state
\end{itemize}

\section*{Indications for Evaluation}
Seek care if:
\begin{itemize}
  \item Severe or worsening symptoms
  \item Symptoms lasting more than acute onset rapid fluctuation requires immediate investigation identifying treatable underlying cause emergent evaluation essential
  \item Accompanied by other concerning signs
\end{itemize}

\section*{Related Symptoms}
\begin{itemize}
  \item Fever
  \item Agitation
  \item Hallucinations
  \item Disorganized thinking
  \item Sleep-wake cycle disruption
\end{itemize}
\textbf{Note:} Always consider serious underlying conditions when evaluating persistent delirium.


\section*{Self-Care / Home Management}
\begin{itemize}
  \item Ensure a safe environment to prevent injury during episodes (remove hazards, avoid driving or operating machinery).
  \item Keep a symptom diary documenting triggers, duration, frequency, and associated features.
  \item Practice stress reduction techniques (deep breathing, meditation, adequate sleep) as stress often exacerbates neurological symptoms.
  \item Avoid alcohol, caffeine, and recreational drugs which can worsen many neurological conditions.
  \item Seek immediate evaluation for sudden-onset, severe, or progressive neurological symptoms.
\end{itemize}

\section*{Diagnostic Approach}
\begin{itemize}
  \item Detailed neurological history includes onset pattern, progression, triggers, associated symptoms, and functional impact.
  \item Comprehensive neurological examination assesses mental status, cranial nerves, motor/sensory function, reflexes, coordination, and gait.
  \item Neuroimaging (CT, MRI) is often indicated to evaluate for structural, vascular, or demyelinating causes.
  \item Specialized testing (EEG for seizures, nerve conduction studies for neuropathy, lumbar puncture for meningitis) is guided by clinical suspicion.
\end{itemize}

\section*{Treatment / Management Overview}
\begin{itemize}
  \item Treatment is directed at the underlying cause identified through diagnostic evaluation.
  \item Symptomatic pharmacotherapy may include analgesics, anticonvulsants, antidepressants, or disease-modifying agents as appropriate.
  \item Physical, occupational, and speech therapy play important roles in functional recovery and rehabilitation.
  \item Referral to neurology is indicated for unexplained, progressive, or treatment-refractory neurological symptoms.
\end{itemize}

\clearpage
